\section{Global coefficients, diagonal characters, and Frobenius precision}\label{sec:coefficients}
The natural class has a fixed global expansion
\begin{equation}\label{eq:globalalpha}
\alpha=a_\phi[\phi]+a_3[3]+a_5[\sqrt5],\qquad a_3a_5\ne0.
\end{equation}
This lies in a Kummer group tensored with $k$. It does not mean that an element of $\F_{49}$ is an exponent of a base field unit. The three coefficients are normalized together, up to one common scalar. Using a different scalar at each completion would discard the global information that makes the reciprocity tests effective.

The actual matrices have the form
\[
\rho(g)=\begin{pmatrix}a(g)&b(g)\\7^nc(g)&d(g)\end{pmatrix}.
\]
Modulo $7^n$, the diagonal entries are characters. Determinant $\chi$ lets us write them as $\chi\eta_n$ and $\eta_n^{-1}$. The natural monodromy restrictions and fixed local traces constrain the same $\eta_n$ at every place. A locally chosen ray character with the right residue does not prove it occurs as this actual diagonal character.

The sequence of Blocks 52 through 58 tightened this coupling. Scalar Frobenius at good thirteen first forced $a_\phi=0$ on that subcase. Later, torsion exponent information at eighty three forced $a_\phi\ne0$ in the retained actual branch. Hence thirteen must be bad there. The conclusion is not a contradiction from two incompatible coefficients. The first coefficient statement assumed good thirteen, so the combined result is $13\mid C$.

At good auxiliary primes, the actual rank two trace obeys a precision condition of the form
\[
\tau\equiv Q+1\pmod{7^r\Oloc}.
\]
The exponent $r$ is supplied by the depth and local diagonal estimates. It is not inferred merely from a rank four congruence. This stronger comparison yields a particularly useful universal lemma.

\begin{lemma}[Discriminant collapse]\label{lem:collapse}
Suppose a good Frobenius polynomial has integral coefficients and factorization
\[
P(T)=(T^2-\tau T+Q)(T^2-\sigma(\tau)T+Q)
 =T^4-ST^3+C_2T^2-QST+Q^2,
\]
where $\Oloc/\Z_7$ is unramified quadratic, $\sigma$ is its nontrivial automorphism, all complex roots have absolute value $\sqrt Q$, and $\tau\equiv Q+1\pmod{7^r\Oloc}$. Then the integer
\[
D=S^2-4(C_2-2Q)
\]
satisfies $7^{2r}\mid D$ and $0\le D\le16Q$. If $7^{2r}>16Q$, then
\[
P(T)=(T^2-aT+Q)^2,
\quad a\in\Z,\quad |a|\le2\sqrt Q,
\quad a\equiv Q+1\pmod{7^r}.
\]
\end{lemma}
\begin{proof}
Expansion gives $S=\tau+\sigma(\tau)$, $C_2-2Q=\tau\sigma(\tau)$, and $D=(\tau-\sigma(\tau))^2$. The congruence is preserved by $\sigma$, so this difference is divisible by $7^r$. Unramifiedness makes its square give ordinary integer divisibility by $7^{2r}$.

The two traces are algebraic roots of $X^2-SX+(C_2-2Q)$. Under their complex embeddings, the corresponding quadratic factors have roots of product $Q$ and magnitude $\sqrt Q$. Their sums are real and lie in $[-2\sqrt Q,2\sqrt Q]$. Their difference therefore has magnitude at most $4\sqrt Q$, proving the bound. If the divisibility exceeds the bound, $D=0$. The common trace is rational and integral, hence an integer $a$, with the stated congruence and magnitude.
\end{proof}
For $Q=2549$ and $r=3$, one has $7^6=117649>40784=16Q$. The allowed integral trace would lie in $[-100,100]$ and be congruent to $149$ modulo $343$. Neither nearest representative, $149$ or $-194$, lies in the interval. Therefore the inherited precision theorem forces $2549\mid C$ whenever it applies at depth at least three.

The rank four congruence alone is weaker. The retained negative control
\[
P(T)=T^4+45T^3+2603T^2+114705T+6497401
\]
is a generic reciprocal Weil polynomial at $2549$ satisfying $(T-1)^2(T-Q)^2$ modulo $343$. Its trace discriminant is $5\cdot49^2$, not divisible by $7^6$. It does not satisfy the actual rank two precision condition. It is not claimed to be a Frey reduction. Its role is to falsify the shortcut that would replace the stronger hypothesis by an unlabelled quartic congruence.

Other fixed primes forced into $C$ at stated depths include the good reduction tests at $41,59,197,491$. Combining these bad primes with reciprocity rows eventually trivialized the diagonal character and reduced depth on declared finite high support classes. The finite support hypotheses remain essential: a new conductor outside the declared universe changes the global ray character presentation.
